\section{Robustness to the substance specification}
\label{sec:robustness}

\S\ref{sec:original} established that the substance component of \citet{lagasio2024} is
underdetermined: the aggregation from per-term TF-IDF values to a document-level $\SUS_i$ is
never stated, and the \texttt{norm} setting is never mentioned. More than one implementation
is therefore admissible on the published description. Of these two silences, this section
takes up the second: the three specifications defined in \S\ref{sec:index} span the
admissible choices of term weighting and vector normalisation. The aggregation silence is
disposed of analytically rather than by comparison, in \S\ref{sec:measurement}. This section
measures how far the index and its labels move across the choices it does span.

Two things this section does not do. It does not choose a specification: the selection
argument rests on agreement with an independently constructed yardstick and is made in
\S\ref{sec:validity}. And it does not treat agreement between two specifications as evidence
that either is correct --- two estimators can agree closely and both be measuring the wrong
quantity, which is precisely the diagnosis \S\ref{sec:measurement} reaches for one of the
three. Throughout, the \emph{score} normalisation is held fixed at $z$-scores; the
sensitivity of the index to that separate choice is the object of \S\ref{sec:artefact}.

\subsection{Design}

The three specifications cross the two open choices one factor at a time. $\susdens$ and
$\sustfidf$ share a denominator --- document length --- and differ only in whether terms are
weighted uniformly or by IDF, so the pair isolates \emph{term weighting}. $\sustfidf$ and
$\suslag$ share IDF weights and differ only in whether the term vector is divided by document
length or scaled to unit $L_2$ norm, so that pair isolates \emph{vector normalisation}.
$\suslag$ implements the inference --- conceded as reasonable in \S\ref{sec:original} and
\S\ref{sec:measurement} --- that the four disclosed \texttt{TfidfVectorizer} settings were the
only departures from library defaults, leaving the default $L_2$ setting in force; the
original never states it in terms. We also report
$\BREADTH$, the effective number of
distinct vocabulary terms (\S\ref{sec:index}): a diagnostic rather than an index component,
and a measure of topical spread that carries no information about how much ESG content a
document holds.

\subsection{Component-level sensitivity}

Table~\ref{tab:robustness-components} gives the correlation matrix over the three
specifications and $\BREADTH$.

\begin{table}[htbp]
\centering
\caption{Pearson correlations among the three substance specifications and $\BREADTH$
($n = 343$).}
\label{tab:robustness-components}
\begin{tabular}{lrrrr}
\toprule
 & $\susdens$ & $\sustfidf$ & $\suslag$ & $\BREADTH$ \\
\midrule
$\susdens$  & 1.0000 & 0.9926 & 0.4177 & 0.4626 \\
$\sustfidf$ & 0.9926 & 1.0000 & 0.4600 & 0.5118 \\
$\suslag$   & 0.4177 & 0.4600 & 1.0000 & 0.9733 \\
$\BREADTH$  & 0.4626 & 0.5118 & 0.9733 & 1.0000 \\
\bottomrule
\end{tabular}
\end{table}

The two contrasts can be read off directly. Across the weighting contrast, $r = 0.9926$:
replacing uniform weights with IDF weights leaves the measurement essentially unchanged.
Across the vector-normalisation contrast, $r = 0.4600$; between the two specifications that differ in
both factors, $r = 0.4177$.

The $\BREADTH$ column identifies what the substitution does. $\suslag$ correlates $0.9733$
with a pure diversity statistic --- higher than it correlates with either of the two
alternative substance measures, and higher than either of those correlates with $\BREADTH$.
Under $L_2$ vector normalisation the score is closer to being a restatement of topical spread than
it is to being a measure of substance. \S\ref{sec:measurement} supplies the mechanism: unit
scaling discards vector magnitude and retains only direction, so what survives is the mix
across the vocabulary, not the amount.

\subsection{Index-level sensitivity and classification agreement}

Table~\ref{tab:robustness-index} carries the comparison through to the assembled index and
its labels.

\begin{table}[htbp]
\centering
\caption{Index correlation and classification agreement against $\ESGSI(\susdens)$
($n = 343$). Labels follow the original's rule, $\ESGSI > 0 \Rightarrow$ ``Potential
ESG-washing''. Score normalisation is $z$-score in every row.}
\label{tab:robustness-index}
\begin{tabular}{llrrr}
\toprule
Substance specification & Factors varied & Pearson $r$ & Flagged & Label flips \\
\midrule
$\susdens$ (reference) & ---                        & 1.0000 & 171/343 & --- \\
$\sustfidf$            & weighting                  & 0.9961 & 167/343 & 6 (1.7\,\%) \\
$\suslag$              & weighting, vector norm.    & 0.6964 & 159/343 & 96 (28.0\,\%) \\
\bottomrule
\end{tabular}
\end{table}

$\suslag$ differs from the reference in both factors, so its 96 flips are not, by themselves,
evidence about which factor drives them; the attribution needs an argument, not a reading of
the table. The argument is a bound, and it uses only the two counts already reported. At most
6 documents anywhere in the corpus disagree under weighting alone (the
$\susdens$--$\sustfidf$ row). Of the 96 documents on which $\susdens$ and $\suslag$ disagree,
therefore, at most 6 can be documents that already disagreed under weighting before vector
normalisation was even applied. At least $96 - 6 = 90$ of the 96 are documents on which
weighting left the label unchanged and vector normalisation alone flipped it --- for those
documents, vector normalisation is the entire explanation, with weighting contributing
nothing. Vector normalisation accounts for at least 90 of the 96 combined label changes;
weighting cannot be credited with more than the 6 flips it produces on its own, anywhere in
the corpus. The ranking in Table~\ref{tab:robustness-index} is not an artefact of comparing
every specification to $\susdens$ rather than to one another: vector normalisation dominates
under any attribution consistent with the two counts already in the table.

Separately, every row of Table~\ref{tab:robustness-index} shares the same $\Z(\SEN)$ term,
and $\SEN$ is only weakly related to any of the three substance measures. Subtracting a
common component that carries little of its own relationship to the compared variables tends
to inflate the correlation between two differences relative to the correlation between the
terms being subtracted --- not as a matter of mathematical necessity, but empirically, and it
does so here: every index-level correlation in Table~\ref{tab:robustness-index} runs above
its component-level counterpart in Table~\ref{tab:robustness-components} (0.9961 against
0.9926; 0.6964 against 0.4177). Read the index-level number accordingly: it is not a tighter
estimate of how closely the underlying substance measurements agree than the component-level
number is. If anything it is the looser of the two, inflated by a term the compared
specifications do not even differ on.

\subsection{Interpretation}

The result is one sentence: \emph{term weighting is second-order and vector normalisation is
first-order}. Substituting uniform weights for IDF moves six documents out of 343; changing
the vector normalisation moves ninety-six. Roughly a quarter of all classifications depend on
a parameter the original never states in terms. What that establishes is a matter of stakes
rather than of bookkeeping, and it should be read alongside the concession
\S\ref{sec:measurement} makes: the \texttt{norm} setting is recoverable from the convention
that only non-default arguments are reported, and we rely on that convention ourselves, so we
do not present the silence as a reproducibility failure. We present it as a first-order choice
left to an inference. The two implementations the published text does not distinguish
\emph{explicitly} disagree about the label of more than one report in four, which is why we
report both rather than settle the question by convention.

The aggregation rule is the opposite case in both respects, and \S\ref{sec:measurement}
disposes of it by argument rather than by comparison: no convention supplies a default for it,
yet because the index is $z$-scored a positive affine rescaling such as switching from a sum
to a mean over the vocabulary cancels out, so that silence does not move a single label even
though nothing in the original resolves it.

It follows that the near-identity of $\susdens$ and $\sustfidf$ is not a validation of either.
It establishes insensitivity to one parameter, which is a property of the parameter, not a
warrant for the measure. Both could be poor measures of substance and would still agree at
$0.9926$. What separates the specifications is external evidence, and that is
\S\ref{sec:validity}'s subject; the present section only quantifies the cost of getting the
choice wrong.

\subsection{Robustness to the vocabulary}
\label{sec:robustness-lexicon}

The section so far varies how the vocabulary is weighted and how the resulting vector is
normalised, never which terms it contains. That leaves untouched the standing objection to any
dictionary-based measure: the result is an artefact of the word list, and a different list would
have produced a different index. \S\ref{sec:lexicon} puts that objection in measured form --- an
external audit of 880 zones from four FY2023 reports identified 147 ESG terms present in the
text and absent from our 154 entries --- and those terms are the material for a robustness check
on the vocabulary itself.

Each of the 143 distinct candidate strings behind that list was measured across all 343 reports
for frequency and dispersion and classified on those statistics (\S\ref{sec:lexicon}); admitting
the dispersed entries and the polysemous ones handled by collocation, and holding back those
flagged sectoral, takes the vocabulary from 154 to 289 lemmatised TF-IDF entries.
Table~\ref{tab:robustness-lexicon} varies it in two steps: the expansion, then the 22 terms held
back as sectoral.

Both experiments run on the same 343 documents and score them with different vocabularies; the
extracted text is held fixed at the expansion's own extraction. That isolation is deliberate and
it is not free: our vocabulary file drives the paragraph-selection regex of
\S\ref{sec:extraction} as well as the TF-IDF vocabulary, so an end-to-end comparison of the two
vocabularies would confound what is scored with what is extracted. Holding the corpus fixed
removes the confound at the cost that the first column of
Table~\ref{tab:robustness-lexicon} is not the index reported elsewhere in this paper: the
adopted specification of \S\ref{sec:lexicon} scores its own extraction, where mean $\susdens$ is
5.897 rather than 5.168. The difference between those two figures is an extraction effect, and
\S\ref{sec:limitations} reports it.

\begin{table}[htbp]
\centering
\caption{Sensitivity of the index to the scoring vocabulary, corpus held fixed ($n = 343$). Each
column after the first is compared with the column immediately to its left. All three columns
score the same extracted text, so the first is the adopted 154-term vocabulary on the
expansion's extraction and not the run reported elsewhere in this paper (see text). Labels
follow the original's rule, $\ESGSI > 0 \Rightarrow$ ``Potential ESG-washing''; score
normalisation is $z$-score in every column. The trend is negative and significant beyond any
conventional level in all three ($p < 10^{-7}$ throughout), with the independence caveat of
\S\ref{sec:temporal} attaching to each.}
\label{tab:robustness-lexicon}
\begin{tabular}{@{}lrrr@{}}
\toprule
 & 154 terms & 289 terms & 289 $+$ 22 sectoral \\
\midrule
Mean $\susdens$                  & 5.168 & 6.516 & 6.541 \\
Change vs preceding column       & ---   & $+26.1$\,\% & $+0.4$\,\% \\
\addlinespace
Flagged                          & 172/343 & 164/343 & 167/343 \\
Label flips vs preceding column  & ---   & 18 (5.2\,\%) & 3 (0.9\,\%) \\
Pearson $r$ vs preceding column  & ---   & 0.9882 & 0.9998 \\
\addlinespace
$\ESGSI(\susdens)$ slope         & $-0.219$/yr & $-0.194$/yr & $-0.193$/yr \\
\bottomrule
\end{tabular}
\end{table}

Near-doubling the vocabulary raises mean substance density by $26.1$\,\%, so the added terms
match text rather than nothing, and it moves the index very little: $r = 0.9882$, 18 of 343
labels change (5.2\,\%), and the flagged count falls from 172/343 to 164/343. The temporal slope
of \S\ref{sec:temporal} goes from $-0.219$ to $-0.194$ per year on this corpus
(\S\ref{sec:temporal-vocabulary}). The 22 terms held back as sectoral move 3 labels of 343 when
added on top of the expansion; \S\ref{sec:limitations} reports what they do at firm level and
why they are excluded.

Two limits bound the check narrowly. The 289-term list is not independent of the corpus it is
tested on: its candidates came from auditing four of these reports and were classified on
statistics computed from all 343. And the check is one-sided, because terms were admitted for
being broadly dispersed, which is precisely the property that makes a term least able to move a
density ranking; a vocabulary built on other principles is untested, and nothing here bounds
what one would do. What the check bounds is the index's sensitivity to a defensible extension of
this vocabulary on this corpus. It does not establish that the vocabulary is right.

The check is not specific to the terms we chose. Deleting 30 of the 154 entries at random and
rescoring --- 5,000 draws --- moves a median of 17 labels at a mean index correlation of 0.988;
our 18 flips and 0.9882 sit at the 56th and 58th percentiles of that distribution, and among the
random draws that displace the index as far as ours does the median is 19 labels against our 18.
What the experiment measures is therefore the index's insensitivity to a vocabulary perturbation
of this magnitude, not the adequacy of the particular terms admitted. Matched on occurrence mass
rather than on term count the curated expansion is modestly the more stable of the two, at the
eighteenth percentile of the random distribution; we report that without leaning on it.

Part of the agreement is structural rather than empirical. The 154 entries are retained whole
inside the 289, so the increment is additive, and its standard deviation is $0.236$ of the base
measure's. That ratio alone bounds the correlation between the two substance scores below at
$\sqrt{1-0.236^{2}} = 0.972$ and the correlation between the two indices at $0.984$. The
observed $0.9882$ lies about a quarter of the way up a window $0.016$ wide. The exercise is not
vacuous --- a fifth of the expanded numerator is new material --- but the range in which the
answer could have landed was fixed by the acceptance rule before any document was scored,
because terms admitted on dispersion contribute a near-constant per-document increment.

Placed beside Table~\ref{tab:robustness-index}, the two experiments complete the set of choices
this paper varies (Table~\ref{tab:robustness-hierarchy}).

\begin{table}[htbp]
\centering
\caption{Labels moved by each of the three choices this paper varies ($n = 343$; $z$-score
normalisation and the $\ESGSI > 0$ threshold throughout). Each row names its own baseline: the
third is an increment on the robustness vocabulary, not on the adopted one. The rows are not
commensurable perturbations and the column does not rank the choices by importance; see the
text.}
\label{tab:robustness-hierarchy}
\begin{tabular}{@{}llr@{}}
\toprule
Choice varied & Measured against & Labels moved \\
\midrule
Vector normalisation (Table~\ref{tab:robustness-index})
  & $\ESGSI(\susdens)$, 154 terms & 96/343 (28.0\,\%) \\
Scoring vocabulary, 154 $\to$ 289 terms
  & $\ESGSI(\susdens)$, 154 terms & 18/343 (5.2\,\%) \\
Sectoral terms, $+22$
  & the 289-term vocabulary      & 3/343 (0.9\,\%) \\
\bottomrule
\end{tabular}
\end{table}

The three rows are not perturbations of comparable size, and the column should not be read as
ranking the choices by how much each matters. Across 2,354 random perturbations of this
vocabulary, label flips scale as $(1-r)^{0.57}$, close to the square root of the index
displacement. On that scaling the three rows differ by a factor of $1.3$ although the
displacements behind them differ by a factor of $1{,}500$: each moves about as many labels as
any perturbation of its size would move. What separates the rows is how far each moved the
index, not what kind of choice each was. Nor is the third row's margin over the second what it
appears: measured against the adopted specification rather than incrementally, the full
$289+22$ vocabulary moves 17 labels, not 3. What the column licenses is narrower than a
hierarchy and is still this section's point: of the three choices we varied, the one
\citet{lagasio2024} never states in terms is the one that changes the most classifications.

The specification adopted for the paper is unchanged --- the 154 entries of \S\ref{sec:lexicon},
on which every figure reported elsewhere is computed. Both vocabularies here are robustness
instruments, not candidates for re-selection: the selection argument is \S\ref{sec:validity}'s,
and it is not re-run on the expanded list.

\subsection{The threshold qualifies the flip counts}

Label flips are counts of documents crossing $\ESGSI = 0$, and that boundary is definitional.
The original states the rule and reports no calibration of it against any external benchmark
(\S\ref{sec:index}). With score normalisation fixed at $z$-score throughout this section, the
consequence is structural: $\ESGSI$ is
a difference of two mean-zero variables, so it has mean zero by construction and the $> 0$
rule flags close to half the corpus under every specification --- 171, 167 and 159 of 343
reports respectively. The cut-off is a property of the sample, not of any firm in it.

Two implications. First, the flip counts are boundary-dependent: imperfectly correlated
indices disagree first about documents lying near the cut-off, so a threshold placed elsewhere
would return different counts. The absolute flip rate is not an error rate and should not be
read as one. Second, the comparison between the two rows is unaffected, because both are
measured against the same boundary on the same sample: whatever the threshold does to $6$, it
does to $96$. The contrast between the two factors is what this section claims; the level is
conditional on a convention.
\S\ref{sec:limitations} takes up the arbitrariness of that convention, and of the relative
ranking it rests on, as a limitation of our own index and not only of the original's.

Finally, the same underdetermination applies to any result computed from this index. The
temporal series in \S\ref{sec:temporal} is therefore reported under all three specifications
rather than under the adopted one alone.
